\section{Circuit Complexity --- Layered DAG}
%%% generated by the blueprint pipeline from TCSlib.Complexity.CircuitComplexity.LayeredDAG
\subsection*{Circuit Complexity --- Layered DAG: Introduction}

Declarations of the Lean module \texttt{TCSlib.Complexity.CircuitComplexity.LayeredDAG}.

\subsection*{Circuit Complexity --- Layered DAG: Results}

\begin{theorem}[Evaluation at the input layer]
\lean{BoolCircuit.LayeredCircuit.evalNode_zero}
\label{BoolCircuit.LayeredCircuit.evalNode_zero}
\leanok
\uses{BoolCircuit.LayeredCircuit.evalNode}
\difficulty{1}
Let $F$ be a layered circuit with input type $\mathrm{inp}$ and alphabet $\alpha$. For a
node $u$ at layer $0$ and an input assignment $xs:\mathrm{inp}\to\alpha$, the value of
$u$ is the value assigned by $xs$ to the input corresponding to $u$.
\end{theorem}

\begin{theorem}[Evaluation at a successor layer]
\lean{BoolCircuit.LayeredCircuit.evalNode_succ}
\label{BoolCircuit.LayeredCircuit.evalNode_succ}
\leanok
\uses{BoolCircuit.LayeredCircuit.Gate.eval, BoolCircuit.LayeredCircuit.evalNode}
\difficulty{1}
Let $F$ be a layered circuit over an alphabet $\alpha$, with input type $\mathrm{inp}$
and output type $\mathrm{out}$. For a natural number $m$ with
$m+1\leq\mathrm{depth}(F)$, a node $u$ at layer $m+1$, and an input assignment
$xs:\mathrm{inp}\to\alpha$, the value of $u$ is the value computed by its gate from the
values of its input nodes at layer $m$.
\end{theorem}

\begin{theorem}[Evaluation with a unique output]
\lean{BoolCircuit.LayeredCircuit.eval₁_eq}
\leanok
\uses{BoolCircuit.LayeredCircuit.evalNode, BoolCircuit.LayeredCircuit.eval₁}
\difficulty{1}
Let $F$ be a layered circuit over $\alpha$ with input type $\mathrm{inp}$ and a unique
output, and let $xs:\mathrm{inp}\to\alpha$ be an input assignment. The circuit's scalar
evaluation equals the value of the node in its last layer corresponding to the unique
output.
\end{theorem}

\begin{definition}[Boolean values as two-element indices]
\lean{BoolCircuit.b2f}
\label{BoolCircuit.b2f}
\leanok
The abbreviation identifies each Boolean value with its corresponding element of the
two-element set $\mathrm{Fin}(2)$, using the inverse of the standard equivalence
between $\mathrm{Fin}(2)$ and the Boolean type.
\end{definition}

\begin{theorem}[Boolean conjunction as multiplication]
\lean{BoolCircuit.b2f_and}
\label{BoolCircuit.b2f_and}
\leanok
\uses{BoolCircuit.b2f}
\difficulty{4}
For Booleans $a$ and $b$, the encoding of their conjunction in $\mathrm{Fin}(2)$ equals
the product of their encodings.
\end{theorem}

\begin{theorem}[Encoding Boolean negation]
\lean{BoolCircuit.b2f_not}
\label{BoolCircuit.b2f_not}
\leanok
\uses{BoolCircuit.b2f}
\difficulty{4}
For every Boolean $a$, the encoding of its negation in $\mathrm{Fin}(2)$ equals $1$
minus the encoding of $a$.
\end{theorem}

\begin{theorem}[Encoding true]
\lean{BoolCircuit.b2f_eq_one_iff}
\label{BoolCircuit.b2f_eq_one_iff}
\leanok
\uses{BoolCircuit.b2f}
\difficulty{4}
For every Boolean $a$, its encoding in $\mathrm{Fin}(2)$ equals $1$ if and only if $a$
is true.
\end{theorem}

\begin{theorem}[Injectivity of the Boolean encoding]
\lean{BoolCircuit.b2f_injective}
\label{BoolCircuit.b2f_injective}
\leanok
\uses{BoolCircuit.b2f}
\difficulty{1}
The encoding of Booleans in $\mathrm{Fin}(2)$ is injective: if two Booleans have the
same encoding, then they are equal.
\end{theorem}

\begin{theorem}[Encoding a conjunction by a product]
\lean{BoolCircuit.prod_b2f_getElem}
\label{BoolCircuit.prod_b2f_getElem}
\leanok
\uses{BoolCircuit.b2f, BoolCircuit.b2f_and}
\difficulty{4}
Encode false as $0$ and true as $1$ in the two-element set with multiplication modulo
$2$. For any finite list $l$ of natural numbers and any Boolean-valued function $f$ on
the natural numbers, the product of the encodings of $f(l_j)$ over all positions $j$ in
the list equals the encoding of the conjunction of the values $f(k)$ for $k$ in $l$.
\end{theorem}

\begin{definition}[Layer gate construction]
\lean{BoolCircuit.layerGateOf}
\label{BoolCircuit.layerGateOf}
\leanok
\uses{BoolCircuit.DAGGate, BoolCircuit.DAGGate.WellFormed, BoolCircuit.Gate}
Let $g$ be a gate whose argument list has no repetitions and whose argument list has
length one whenever $g$ is a NOT gate. For any domain $D$ and any assignment of a vertex
of $D$ to each argument in the list, the construction gives a gate with values in
$	ext{Fin}(2)$, wired to those assigned vertices. Its operation is $x\mapsto 1-x(0)$ for
a NOT gate, and the product of the input values for either an AND or an OR gate.
\end{definition}

\begin{theorem}[Standard operation for a layered gate]
\lean{BoolCircuit.layerGateOf_mem}
\label{BoolCircuit.layerGateOf_mem}
\leanok
\uses{BoolCircuit.DAGGate, BoolCircuit.DAGGate.WellFormed, BoolCircuit.layerGateOf, BoolCircuit.stdGateOps}
\difficulty{4}
Let $g$ be a Boolean DAG gate whose argument list has no repetitions and whose NOT gate,
if it is a NOT gate, has exactly one argument. For any type $D$ and any assignment of an
element of $D$ to each argument occurrence of $g$, the operation of the resulting
layered gate belongs to the standard set consisting of the identity, NOT, and
arbitrary-fanin AND operations on $\mathrm{Fin}\,2$.
\end{theorem}

\begin{theorem}[Evaluation of a layered Boolean gate]
\lean{BoolCircuit.layerGateOf_eval}
\label{BoolCircuit.layerGateOf_eval}
\leanok
\uses{BoolCircuit.DAGGate, BoolCircuit.DAGGate.WellFormed, BoolCircuit.DAGGate.eval, BoolCircuit.LayeredCircuit.Gate.eval, BoolCircuit.b2f, BoolCircuit.b2f_not, BoolCircuit.layerGateOf, BoolCircuit.prod_b2f_getElem}
\difficulty{5}
Let $g$ be a Boolean DAG gate whose argument list has no repetitions, with exactly one
argument if it is a NOT gate, and suppose $g$ is not an OR gate. Let $D$ be a type, let
$e$ assign an element of $D$ to each argument position in the list, let
$v:D\to\mathrm{Fin}\,2$, and let $w$ be a Boolean list. If, for every argument index $a$
of $g$, the value $v(e(a))$ is the $0$–$1$ encoding of the value of $w$ at index $a$
(taking that value to be false when the index is out of range), then the gate on $D$
constructed from $g$ and $e$, evaluated using $v$, equals the $0$–$1$ encoding of the
evaluation of $g$ on $w$.
\end{theorem}

\begin{definition}[Nodes at a circuit layer]
\lean{BoolCircuit.DAGCircuit.layerNodes}
\label{BoolCircuit.DAGCircuit.layerNodes}
\leanok
For a directed acyclic circuit with $n$ input vertices and $L$ gates, the node type at
layer $d$ is the set of indices $\mathrm{Fin}(n+d)$ when $d\le L$, and is a singleton
when $d>L$, where $d$ ranges over the indices from $0$ through $L+1$.
\end{definition}

\begin{definition}[Input indices as layer nodes]
\lean{BoolCircuit.DAGCircuit.toNode}
\label{BoolCircuit.DAGCircuit.toNode}
\leanok
\uses{BoolCircuit.DAGCircuit.layerNodes}
Let $C$ be a directed acyclic Boolean circuit with $n$ inputs and a finite list of
gates. For a layer index $d$ satisfying $d\leq$ the number of gates and an index
$i\in\{0,\ldots,n+d-1\}$, the corresponding node at layer $d$ is the node indexed by
$i$.
\end{definition}

\begin{definition}[Index representation of layer nodes]
\lean{BoolCircuit.DAGCircuit.ofNode}
\label{BoolCircuit.DAGCircuit.ofNode}
\leanok
\uses{BoolCircuit.DAGCircuit.layerNodes}
Let $C$ be a directed acyclic circuit with $n$ input vertices, and let $d$ be a layer
index with $d$ at most the number of gates. Every node $u$ in layer $d$ is represented
by its corresponding index in $	ext{Fin}(n+d)$.
\end{definition}

\begin{theorem}[Inverse conversion between layer indices]
\lean{BoolCircuit.DAGCircuit.ofNode_toNode}
\label{BoolCircuit.DAGCircuit.ofNode_toNode}
\leanok
\uses{BoolCircuit.DAGCircuit.ofNode, BoolCircuit.DAGCircuit.toNode}
\difficulty{1}
Let $C$ be a directed acyclic circuit with $n$ input vertices and a list of gates, and
let $d$ be a layer index with $d\le$ the number of gates. Identifying the vertices in
that layer with the indices $0,\ldots,n+d-1$, converting such an index to a layer vertex
and then back to an index returns the original index.
\end{theorem}

\begin{theorem}[Bound on the cast successor index]
\lean{BoolCircuit.DAGCircuit.castSucc_le}
\label{BoolCircuit.DAGCircuit.castSucc_le}
\leanok
\difficulty{3}
Let $C$ be a directed acyclic Boolean circuit with $n$ inputs and $N$ gates. For every
$d$ with $0\le d\le N$, the value of $d$ regarded as an element of the successor
embedding into indices from $0$ through $N+1$ is at most $N$.
\end{theorem}

\begin{theorem}[Successor bound]
\lean{BoolCircuit.DAGCircuit.succ_le}
\label{BoolCircuit.DAGCircuit.succ_le}
\leanok
\difficulty{2}
Let $m$ be the number of gates in a directed acyclic circuit, and let $d$ be a natural
number with $d<m+1$. If $d<m$, then the successor of $d$ is at most $m$.
\end{theorem}

\begin{definition}[Gate in a circuit layer]
\lean{BoolCircuit.DAGCircuit.layerGate}
\label{BoolCircuit.DAGCircuit.layerGate}
\leanok
\uses{BoolCircuit.DAGCircuit.IsWellFormed, BoolCircuit.DAGCircuit.castSucc_le, BoolCircuit.DAGCircuit.layerNodes, BoolCircuit.DAGCircuit.ofNode, BoolCircuit.DAGCircuit.succ_le, BoolCircuit.DAGCircuit.toNode, BoolCircuit.Gate, BoolCircuit.LayeredCircuit.GateOp.id, BoolCircuit.layerGateOf}
Let $C$ be a directed acyclic Boolean circuit with $n$ input nodes, an ordered list of
$m$ gates, and an output node. Assume each gate has distinct input arguments and every
NOT gate has exactly one argument. For an integer $d$ with $0 \le d \le m$ and a node
$u$ in layer $d+1$, define a gate whose input wires come from layer $d$ as follows: if
$d<m$ and the circuit node represented by $u$ has index less than $n+d$, the gate copies
that node's value; otherwise, it is the gate at position $d$, with its arguments wired
to their corresponding nodes in layer $d$. When $d=m$, the gate copies the circuit's
output node.
\end{definition}

\begin{definition}[Layering a well-formed circuit]
\lean{BoolCircuit.DAGCircuit.toLayered}
\label{BoolCircuit.DAGCircuit.toLayered}
\leanok
\uses{BoolCircuit.DAGCircuit.IsWellFormed, BoolCircuit.DAGCircuit.layerGate, BoolCircuit.DAGCircuit.layerNodes, BoolCircuit.LayeredCircuit}
Let $C$ be a directed acyclic circuit with $n$ inputs, and suppose each gate has
distinct input nodes and every NOT gate has exactly one input. Its associated layered
circuit has depth one more than the number of gates, input type $	ext{Fin}\,n$, and
singleton output type; its layers and gates are those determined by $C$.
\end{definition}

\begin{theorem}[Finiteness of the layered circuit]
\lean{BoolCircuit.DAGCircuit.toLayered_finite}
\label{BoolCircuit.DAGCircuit.toLayered_finite}
\leanok
\uses{BoolCircuit.DAGCircuit.layerNodes, BoolCircuit.DAGCircuit.toLayered, BoolCircuit.LayeredCircuit.Finite}
\difficulty{1}
Let $C$ be a DAG circuit on $n$ inputs whose gates have distinct arguments and whose NOT
gates each have exactly one argument. Every layer of the layered circuit associated to
$C$ has finitely many nodes.
\end{theorem}

\begin{theorem}[Gate operations in the layered circuit]
\lean{BoolCircuit.DAGCircuit.toLayered_onlyUsesGates}
\label{BoolCircuit.DAGCircuit.toLayered_onlyUsesGates}
\leanok
\uses{BoolCircuit.DAGCircuit.layerGate, BoolCircuit.DAGCircuit.toLayered, BoolCircuit.LayeredCircuit.onlyUsesGates, BoolCircuit.layerGateOf_mem, BoolCircuit.stdGateOps}
\difficulty{4}
Let $C$ be a DAG circuit on $n$ inputs whose gates have distinct arguments and whose NOT
gates each have exactly one argument. Every gate operation in the associated layered
circuit is either the identity, the NOT operation, or an unbounded fan-in AND operation.
\end{theorem}

\begin{theorem}[Evaluation of layered DAG nodes]
\lean{BoolCircuit.DAGCircuit.toLayered_evalNode}
\label{BoolCircuit.DAGCircuit.toLayered_evalNode}
\leanok
\uses{BoolCircuit.DAGCircuit.layerGate, BoolCircuit.DAGCircuit.layerNodes, BoolCircuit.DAGCircuit.ofNode, BoolCircuit.DAGCircuit.ofNode_toNode, BoolCircuit.DAGCircuit.succ_le, BoolCircuit.DAGCircuit.toLayered, BoolCircuit.DAGCircuit.values, BoolCircuit.DAGCircuit.values_getD_gate, BoolCircuit.DAGCircuit.values_getD_input, BoolCircuit.LayeredCircuit.Gate.eval, BoolCircuit.LayeredCircuit.evalNode, BoolCircuit.LayeredCircuit.evalNode_succ, BoolCircuit.LayeredCircuit.evalNode_zero, BoolCircuit.b2f, BoolCircuit.layerGateOf_eval}
\difficulty{6}
Let $C$ be a Boolean DAG circuit on $n$ inputs whose gates have distinct arguments and
whose NOT gates each have exactly one argument. Suppose no gate of $C$ is an OR gate,
and let $x$ be a Boolean input assignment. For every $m$ with $0\le m\le
|C.\mathrm{gates}|$ and every node $u$ in layer $m$, the value of $u$ in the layered
circuit associated to $C$ is the Boolean value, encoded in $\mathrm{Fin}\,2$, of the
corresponding node in the execution of $C$; a node index outside the execution’s value
list is assigned false.
\end{theorem}

\begin{theorem}[Evaluation of the layered circuit]
\lean{BoolCircuit.DAGCircuit.toLayered_eval₁}
\leanok
\uses{BoolCircuit.DAGCircuit.eval, BoolCircuit.DAGCircuit.layerGate, BoolCircuit.DAGCircuit.layerNodes, BoolCircuit.DAGCircuit.ofNode_toNode, BoolCircuit.DAGCircuit.toLayered, BoolCircuit.DAGCircuit.toLayered_evalNode, BoolCircuit.LayeredCircuit.Gate.eval, BoolCircuit.LayeredCircuit.evalNode, BoolCircuit.LayeredCircuit.evalNode_succ, BoolCircuit.LayeredCircuit.eval₁, BoolCircuit.b2f}
\difficulty{4}
Let $C$ be a Boolean DAG circuit on $n$ inputs whose gates have distinct arguments and
whose NOT gates each have exactly one argument. Suppose no gate of $C$ is an OR gate.
For every Boolean input assignment $x$, the value of the associated layered circuit,
whose output type has a unique element, equals the output value of $C$, with Boolean
values encoded in $\mathrm{Fin}\,2$.
\end{theorem}

\begin{theorem}[Size bound for the layered circuit]
\lean{BoolCircuit.DAGCircuit.toLayered_size_le}
\label{BoolCircuit.DAGCircuit.toLayered_size_le}
\leanok
\uses{BoolCircuit.DAGCircuit.layerNodes, BoolCircuit.DAGCircuit.size, BoolCircuit.DAGCircuit.toLayered, BoolCircuit.LayeredCircuit.size}
\difficulty{5}
Let $C$ be a DAG circuit on $n$ inputs whose gates have distinct arguments and whose NOT
gates each have exactly one argument. The size of the associated layered circuit is at
most $(|C.\mathrm{gates}|+1)(\mathrm{size}(C)+1)$, where
$\mathrm{size}(C)=n+|C.\mathrm{gates}|$.
\end{theorem}

\begin{definition}[Boolean gate shapes]
\lean{BoolCircuit.GateShape}
\label{BoolCircuit.GateShape}
\leanok
For a type $D$ of input nodes, a gate shape is either a wire from one node $a\in D$, a
negation of one node $a\in D$, or a conjunction of a finite list of nodes from $D$.
\end{definition}

\begin{definition}[Evaluation of a gate shape]
\lean{BoolCircuit.GateShape.eval}
\label{BoolCircuit.GateShape.eval}
\leanok
\uses{BoolCircuit.GateShape}
Given Boolean values on the nodes in a type $D$, the value of a wire is the value at its
node, the value of a negation is one minus the value at its node, and the value of a
conjunction is the product of the values at the nodes in its list.
\end{definition}

\begin{theorem}[Standard gate operations have simple shapes]
\lean{BoolCircuit.exists_gateShape}
\label{BoolCircuit.exists_gateShape}
\leanok
\uses{BoolCircuit.Gate, BoolCircuit.GateShape, BoolCircuit.GateShape.eval, BoolCircuit.LayeredCircuit.Gate.eval, BoolCircuit.stdGateOps}
\difficulty{4}
Let $D$ be a type, and let $g$ be a gate with Boolean outputs whose operation is either
the identity operation, a unary negation, or a conjunction of some arity. Then there is
a shape built from a wire, a negated wire, or a conjunction of a list of elements of $D$
whose value agrees with the value of $g$ for every assignment of Boolean values to $D$.
\end{theorem}

\begin{definition}[Shape of a standard gate]
\lean{BoolCircuit.gateShape}
\label{BoolCircuit.gateShape}
\leanok
\uses{BoolCircuit.Gate, BoolCircuit.GateShape, BoolCircuit.exists_gateShape, BoolCircuit.stdGateOps}
For a gate on Boolean values whose operation is either the identity, NOT, or an
unbounded fan-in AND, choose a gate shape representing that operation with the gate's
inputs wired to nodes of its input domain.
\end{definition}

\begin{theorem}[Boolean gate representation by a shape]
\lean{BoolCircuit.gateShape_spec}
\label{BoolCircuit.gateShape_spec}
\leanok
\uses{BoolCircuit.Gate, BoolCircuit.GateShape.eval, BoolCircuit.LayeredCircuit.Gate.eval, BoolCircuit.exists_gateShape, BoolCircuit.gateShape, BoolCircuit.stdGateOps}
\difficulty{1}
Let $D$ be a set, and let $g$ be a gate with input domain $D$ whose operation is either
the identity, negation, or an AND operation of some arity. There is a gate shape built
from a wire, a negated wire, or a conjunction of wires such that, for every assignment
of values in $\mathrm{Fin}\,2$ to $D$, evaluating $g$ gives the same value as
evaluating that shape.
\end{theorem}

\begin{theorem}[Deduplication preserves a list-wide test]
\lean{BoolCircuit.all_dedup'}
\label{BoolCircuit.all_dedup'}
\leanok
\difficulty{3}
For every list $l$ of natural numbers and every Boolean-valued function $p$ on the
natural numbers, applying the test that every entry satisfies $p$ gives the same result
on $l$ and on the list obtained by removing repeated entries from $l$.
\end{theorem}

\begin{theorem}[Boolean conjunction as a product]
\lean{BoolCircuit.b2f_all}
\label{BoolCircuit.b2f_all}
\leanok
\uses{BoolCircuit.b2f, BoolCircuit.b2f_and}
\difficulty{4}
For any type $\alpha$, list $l : \mathrm{List}(\alpha)$, and function $f : \alpha \to
\{\mathrm{false},\mathrm{true}\}$, encode each Boolean by its corresponding element of
$\mathrm{Fin}(2)$. The encoding of the conjunction of the values $f(x)$ over $x$ in $l$
equals the product in $\mathrm{Fin}(2)$ of their encodings.
\end{theorem}

\begin{definition}[Size of a circuit layer]
\lean{BoolCircuit.LayeredCircuit.succCard}
\label{BoolCircuit.LayeredCircuit.succCard}
\leanok
Let $F$ be a finite layered Boolean circuit from $	ext{Fin}\,n$ to a singleton, using
only identity, NOT, and unbounded fan-in AND gates. For each natural number $j$, define
its successor-layer size to be the number of nodes in layer $j+1$ when $j$ is less than
the circuit depth, and zero otherwise.
\end{definition}

\begin{definition}[Cumulative layer offset]
\lean{BoolCircuit.LayeredCircuit.offset}
\label{BoolCircuit.LayeredCircuit.offset}
\leanok
\uses{BoolCircuit.LayeredCircuit.succCard}
Let $F$ be a finite layered Boolean circuit from $	ext{Fin}\,n$ to a singleton, using
only identity, NOT, and unbounded fan-in AND gates. For each natural number $m$, define
its offset to be the sum of the successor-layer sizes over all layers indexed by $j<m$.
\end{definition}

\begin{theorem}[Successor formula for the layer offset]
\lean{BoolCircuit.LayeredCircuit.offset_succ}
\label{BoolCircuit.LayeredCircuit.offset_succ}
\leanok
\uses{BoolCircuit.LayeredCircuit.offset, BoolCircuit.LayeredCircuit.succCard}
\difficulty{1}
Let $F$ be a finite layered Boolean circuit of depth $d$ using only the standard gate
operations, and let its offset at $m$ be the sum of the numbers of nodes in layers $j+1$
for $j<m$, counting zero when $j\ge d$. Then for every natural number $m$, the offset at
$m+1$ is the offset at $m$ plus the number of nodes in layer $m+1$ if $m<d$, and is just
the offset at $m$ otherwise.
\end{theorem}

\begin{definition}[Enumeration of a circuit layer]
\lean{BoolCircuit.LayeredCircuit.layerEnum}
\label{BoolCircuit.LayeredCircuit.layerEnum}
\leanok
Let $F$ be a finite layered circuit with Boolean-valued gates, input set $	ext{Fin}\,n$,
and a singleton output, using only the standard gate operations. For every natural
number $m$ with $m < 	ext{depth}(F)+1$, the nodes in layer $m$ are in bijection with the
natural numbers less than the number of nodes in that layer.
\end{definition}

\begin{definition}[Natural-number labels for circuit vertices]
\lean{BoolCircuit.LayeredCircuit.layerVertex}
\label{BoolCircuit.LayeredCircuit.layerVertex}
\leanok
\uses{BoolCircuit.LayeredCircuit.layerEnum, BoolCircuit.LayeredCircuit.offset}
Let $F$ be a layered circuit with Boolean values and $n$ input wires, and let $m$ be a
nonnegative integer with $m \le \mathrm{depth}(F)+1$. Assign each node $u$ in layer $m$ a
natural-number label as follows: in layer $0$, its label is the index of the
corresponding input wire; in layer $m+1$, its label is $n$ plus the total number of
successor nodes in layers $0$ through $m-1$, plus the position of $u$ in an enumeration
of layer $m+1$.
\end{definition}

\begin{definition}[Gate at a circuit layer]
\lean{BoolCircuit.LayeredCircuit.layerGate}
\label{BoolCircuit.LayeredCircuit.layerGate}
\leanok
\uses{BoolCircuit.DAGGate, BoolCircuit.LayeredCircuit.layerVertex, BoolCircuit.gateShape}
Let $F$ be a finite layered Boolean circuit that uses only identity, negation, and
unbounded fan-in conjunction gates. For a layer $m$ below the circuit depth and a node
$u$ in layer $m+1$, define its associated directed acyclic graph gate by reading the
gate’s shape: an identity wire becomes a one-input AND gate, a negation becomes a NOT
gate, and a conjunction becomes an AND gate whose inputs are the corresponding layer
vertices, with repeated inputs removed.
\end{definition}

\begin{definition}[Gates through a layer bound]
\lean{BoolCircuit.LayeredCircuit.gatesUpTo}
\label{BoolCircuit.LayeredCircuit.gatesUpTo}
\leanok
\uses{BoolCircuit.DAGGate, BoolCircuit.LayeredCircuit.layerEnum, BoolCircuit.LayeredCircuit.layerGate}
Let $F$ be a finite layered Boolean circuit using only identity, negation, and unbounded
fan-in conjunction gates. For each natural number $m$, define a list of directed acyclic
graph gates recursively: the list for $0$ is empty, and the list for $m+1$ appends all
gates from layer $m$ when $m$ is below the circuit depth, or appends nothing otherwise;
each layer’s gates are listed by an enumeration of its nodes.
\end{definition}

\begin{theorem}[Length of the circuit gate list]
\lean{BoolCircuit.LayeredCircuit.length_gatesUpTo}
\label{BoolCircuit.LayeredCircuit.length_gatesUpTo}
\leanok
\uses{BoolCircuit.LayeredCircuit.gatesUpTo, BoolCircuit.LayeredCircuit.offset, BoolCircuit.LayeredCircuit.offset_succ, BoolCircuit.LayeredCircuit.succCard}
\difficulty{4}
Let $F$ be a finite layered Boolean circuit of depth $d$ using only the standard gate
operations. For every natural number $m$, the list of gates in the first $m$ layers has
length equal to the sum, over $j<m$, of the number of nodes in layer $j+1$ when $j<d$,
with each summand taken as zero when $j\ge d$.
\end{theorem}

\begin{theorem}[Monotonicity of the gate lists]
\lean{BoolCircuit.LayeredCircuit.gatesUpTo_mono}
\label{BoolCircuit.LayeredCircuit.gatesUpTo_mono}
\leanok
\uses{BoolCircuit.LayeredCircuit.gatesUpTo}
\difficulty{5}
Let $F$ be a finite layered Boolean circuit whose gates use only identity, NOT, and
unbounded fan-in AND operations. For any natural numbers $m\le M$, the list of gates
through layer $M$ is the list of gates through layer $m$ followed by some list of
additional gates.
\end{theorem}

\begin{theorem}[Upper bound for layer vertex indices]
\lean{BoolCircuit.LayeredCircuit.layerVertex_lt}
\label{BoolCircuit.LayeredCircuit.layerVertex_lt}
\leanok
\uses{BoolCircuit.LayeredCircuit.layerEnum, BoolCircuit.LayeredCircuit.layerVertex, BoolCircuit.LayeredCircuit.offset, BoolCircuit.LayeredCircuit.offset_succ, BoolCircuit.LayeredCircuit.succCard}
\difficulty{4}
Let $F$ be a finite layered Boolean circuit whose gates use only identity, NOT, and
unbounded fan-in AND operations. For every layer $m<\mathrm{depth}(F)+1$ and every node
$u$ in layer $m$, the index assigned to $u$ is strictly less than $n$ plus the
cumulative number of nodes in the successor layers preceding layer $m$.
\end{theorem}

\begin{theorem}[Lower bound for layer vertex indices]
\lean{BoolCircuit.LayeredCircuit.layerVertex_ge}
\label{BoolCircuit.LayeredCircuit.layerVertex_ge}
\leanok
\uses{BoolCircuit.LayeredCircuit.layerVertex, BoolCircuit.LayeredCircuit.offset}
\difficulty{1}
Let $F$ be a finite layered Boolean circuit whose gates use only identity, NOT, and
unbounded fan-in AND operations. If $m+1<\mathrm{depth}(F)+1$ and $u$ is a node in layer
$m+1$, then the index assigned to $u$ is at least $n$ plus the cumulative number of
nodes in the successor layers preceding layer $m$.
\end{theorem}

\begin{theorem}[Layer gate position in the gate list]
\lean{BoolCircuit.LayeredCircuit.gatesUpTo_getElem}
\label{BoolCircuit.LayeredCircuit.gatesUpTo_getElem}
\leanok
\uses{BoolCircuit.LayeredCircuit.gatesUpTo, BoolCircuit.LayeredCircuit.layerEnum, BoolCircuit.LayeredCircuit.layerGate, BoolCircuit.LayeredCircuit.length_gatesUpTo, BoolCircuit.LayeredCircuit.offset}
\difficulty{3}
Let $F$ be a finite layered Boolean circuit whose gates use only identity, NOT, and
unbounded fan-in AND operations. For a layer $m$ below the circuit depth and a node $u$
in the successor layer, if the index given by the cumulative number of nodes in earlier
successor layers plus the rank of $u$ in its layer is within the list of gates through
layer $m$, then the gate at that index is precisely the gate associated with $u$ in
layer $m$.
\end{theorem}

\begin{theorem}[Layer gate arguments precede its layer]
\lean{BoolCircuit.LayeredCircuit.layerGate_args_lt}
\label{BoolCircuit.LayeredCircuit.layerGate_args_lt}
\leanok
\uses{BoolCircuit.LayeredCircuit.layerGate, BoolCircuit.LayeredCircuit.layerVertex_lt, BoolCircuit.LayeredCircuit.offset}
\difficulty{4}
Let $F$ be a finite layered Boolean circuit whose gates use only identity, NOT, and
unbounded fan-in AND operations. If $m$ is below the depth of $F$ and $u$ is a node in
the successor of layer $m$, then every argument vertex of the gate associated with $u$
in layer $m$ has index strictly less than $n$ plus the cumulative number of nodes in the
successor layers preceding layer $m$.
\end{theorem}

\begin{theorem}[Layer gates are well formed]
\lean{BoolCircuit.LayeredCircuit.layerGate_wellFormed}
\label{BoolCircuit.LayeredCircuit.layerGate_wellFormed}
\leanok
\uses{BoolCircuit.DAGGate.WellFormed, BoolCircuit.LayeredCircuit.layerGate}
\difficulty{4}
Let $F$ be a finite layered Boolean circuit whose gates use only identity, NOT, and
unbounded fan-in AND operations. For every layer $m$ below the depth of $F$ and every
node $u$ in the successor of that layer, the associated gate has pairwise distinct
argument vertices, and if its kind is NOT then it has exactly one argument.
\end{theorem}

\begin{theorem}[Acyclicity of the bounded gate list]
\lean{BoolCircuit.LayeredCircuit.gatesUpTo_acyclic}
\label{BoolCircuit.LayeredCircuit.gatesUpTo_acyclic}
\leanok
\uses{BoolCircuit.GatesAcyclic, BoolCircuit.GatesAcyclic.nil, BoolCircuit.LayeredCircuit.gatesUpTo, BoolCircuit.LayeredCircuit.layerGate_args_lt, BoolCircuit.LayeredCircuit.length_gatesUpTo}
\difficulty{5}
Let $F$ be a layered Boolean circuit with $n$ inputs and a unique output. Suppose every
layer has finitely many nodes and every gate operation is either the identity, NOT, or
an unbounded fan-in AND. For every natural number $m$, the list formed from the first
$m$ layers of gates is acyclic: every source index of the gate at position $i$ is less
than $n+i$.
\end{theorem}

\begin{theorem}[Layer vertex values match circuit evaluation]
\lean{BoolCircuit.LayeredCircuit.layerVertex_value}
\label{BoolCircuit.LayeredCircuit.layerVertex_value}
\leanok
\uses{BoolCircuit.DAGGate.eval, BoolCircuit.DAGGate.eval_congr, BoolCircuit.GateShape.eval, BoolCircuit.LayeredCircuit.evalNode, BoolCircuit.LayeredCircuit.evalNode_succ, BoolCircuit.LayeredCircuit.evalNode_zero, BoolCircuit.LayeredCircuit.gatesUpTo, BoolCircuit.LayeredCircuit.gatesUpTo_getElem, BoolCircuit.LayeredCircuit.gatesUpTo_mono, BoolCircuit.LayeredCircuit.layerEnum, BoolCircuit.LayeredCircuit.layerGate, BoolCircuit.LayeredCircuit.layerGate_args_lt, BoolCircuit.LayeredCircuit.layerVertex, BoolCircuit.LayeredCircuit.length_gatesUpTo, BoolCircuit.LayeredCircuit.offset, BoolCircuit.LayeredCircuit.offset_succ, BoolCircuit.LayeredCircuit.succCard, BoolCircuit.all_dedup', BoolCircuit.b2f, BoolCircuit.b2f_all, BoolCircuit.b2f_not, BoolCircuit.gateShape_spec, BoolCircuit.runWith, BoolCircuit.runWith_getD_gate, BoolCircuit.runWith_getD_take, BoolCircuit.vertexValue, BoolCircuit.vertexValue_append, BoolCircuit.vertexValue_input}
\difficulty{6}
Let $F$ be a finite layered Boolean circuit whose gates use only identity, NOT, and
unbounded fan-in AND operations, and let $x$ assign a Boolean value to each of its $n$
inputs. For every layer $m<\mathrm{depth}(F)+1$ and every node $u$ in that layer, the
Boolean value produced by the gate list through layer $m$ at the index assigned to $u$
equals the value of $u$ when the layered circuit is evaluated on $x$.
\end{theorem}

\begin{definition}[Directed acyclic graph of a layered circuit]
\lean{BoolCircuit.LayeredCircuit.toDAG}
\label{BoolCircuit.LayeredCircuit.toDAG}
\leanok
\uses{BoolCircuit.DAGCircuit, BoolCircuit.LayeredCircuit.gatesUpTo, BoolCircuit.LayeredCircuit.gatesUpTo_acyclic, BoolCircuit.LayeredCircuit.layerVertex, BoolCircuit.LayeredCircuit.layerVertex_lt, BoolCircuit.LayeredCircuit.length_gatesUpTo}
Let $F$ be a finite layered Boolean circuit with $n$ input bits that uses only identity,
negation, and unbounded fan-in conjunction gates. Its associated directed acyclic graph
circuit has the gates obtained from all layers in order, with each gate reading only
earlier vertices, and has as output the vertex corresponding to the circuit’s final
output.
\end{definition}

\begin{theorem}[Evaluation preserved by conversion to a DAG]
\lean{BoolCircuit.LayeredCircuit.toDAG_eval}
\label{BoolCircuit.LayeredCircuit.toDAG_eval}
\leanok
\uses{BoolCircuit.DAGCircuit.eval, BoolCircuit.LayeredCircuit.eval₁, BoolCircuit.LayeredCircuit.layerVertex_value, BoolCircuit.LayeredCircuit.toDAG, BoolCircuit.b2f}
\difficulty{1}
Let $F$ be a finite layered Boolean circuit with $n$ inputs that uses only the standard
gate operations, and let $x$ be an assignment of Boolean values to its inputs. The
Boolean value computed by the directed acyclic circuit obtained from $F$, viewed as an
element of $\mathrm{Fin}\,2$, equals the value computed by $F$ at its unique output,
under the input assignment obtained by encoding each Boolean value as its corresponding
element of $\mathrm{Fin}\,2$.
\end{theorem}

\begin{theorem}[Size of the converted circuit]
\lean{BoolCircuit.LayeredCircuit.toDAG_size}
\label{BoolCircuit.LayeredCircuit.toDAG_size}
\leanok
\uses{BoolCircuit.DAGCircuit.size, BoolCircuit.LayeredCircuit.length_gatesUpTo, BoolCircuit.LayeredCircuit.offset, BoolCircuit.LayeredCircuit.size, BoolCircuit.LayeredCircuit.succCard, BoolCircuit.LayeredCircuit.toDAG}
\difficulty{4}
Let $F$ be a finite layered Boolean circuit with $n$ inputs that uses only the standard
gate operations. The size of the directed acyclic circuit obtained from $F$ is $n$ plus
the size of $F$.
\end{theorem}

\begin{theorem}[Well-formed gates in each prefix]
\lean{BoolCircuit.LayeredCircuit.gatesUpTo_wellFormed}
\label{BoolCircuit.LayeredCircuit.gatesUpTo_wellFormed}
\leanok
\uses{BoolCircuit.DAGGate.WellFormed, BoolCircuit.LayeredCircuit.gatesUpTo, BoolCircuit.LayeredCircuit.layerGate_wellFormed}
\difficulty{4}
Let $F$ be a finite layered Boolean circuit whose gates use only identity, NOT, and
unbounded fan-in AND operations. For every natural number $M$, each gate in the list
through layer $M$ has pairwise distinct argument vertices, and if its kind is NOT then
it has exactly one argument.
\end{theorem}

\begin{theorem}[Well-formedness of the converted circuit]
\lean{BoolCircuit.LayeredCircuit.toDAG_isWellFormed}
\label{BoolCircuit.LayeredCircuit.toDAG_isWellFormed}
\leanok
\uses{BoolCircuit.DAGCircuit.IsWellFormed, BoolCircuit.LayeredCircuit.gatesUpTo_wellFormed, BoolCircuit.LayeredCircuit.toDAG}
\difficulty{1}
Let $F$ be a finite layered Boolean circuit that uses only the standard gate operations.
The directed acyclic circuit obtained from $F$ is well formed: the argument list of each
gate has no repetitions, and every NOT gate has exactly one argument.
\end{theorem}

\begin{definition}[Layered circuit associated to a tree circuit]
\lean{BoolCircuit.TreeCircuit.toLayered}
\label{BoolCircuit.TreeCircuit.toLayered}
\leanok
\uses{BoolCircuit.DAGCircuit.deMorgan, BoolCircuit.DAGCircuit.deMorgan_isWellFormed, BoolCircuit.DAGCircuit.toLayered, BoolCircuit.LayeredCircuit, BoolCircuit.TreeCircuit.toDAG, BoolCircuit.TreeCircuit.toDAG_isWellFormed}
Let $c$ be a tree Boolean circuit with $n$ input bits. Define its associated layered
Boolean circuit by first converting it to a directed acyclic graph circuit using De
Morgan transformations, then converting that graph circuit to a layered circuit with the
same input bits and a single output.
\end{definition}

\begin{theorem}[Evaluation of the layered circuit]
\lean{BoolCircuit.TreeCircuit.toLayered_eval₁}
\leanok
\uses{BoolCircuit.DAGCircuit.deMorgan_eval, BoolCircuit.DAGCircuit.deMorgan_kind_ne_or, BoolCircuit.DAGCircuit.toLayered_eval₁, BoolCircuit.LayeredCircuit.eval₁, BoolCircuit.TreeCircuit.eval, BoolCircuit.TreeCircuit.toDAG_eval, BoolCircuit.TreeCircuit.toLayered, BoolCircuit.b2f}
\difficulty{3}
Let $c$ be a Boolean circuit on $n$ variables and let $x : \mathrm{Fin}\,n \to \bbf_2$
be an input assignment. Evaluating the layered circuit associated to $c$ on the input
obtained by mapping each Boolean value into $\mathrm{Fin}\,2$ gives the value of $c$ on
$x$, mapped into $\mathrm{Fin}\,2$.
\end{theorem}

\begin{theorem}[Finiteness of the circuit layers]
\lean{BoolCircuit.TreeCircuit.toLayered_finite}
\label{BoolCircuit.TreeCircuit.toLayered_finite}
\leanok
\uses{BoolCircuit.DAGCircuit.toLayered_finite, BoolCircuit.LayeredCircuit.Finite, BoolCircuit.TreeCircuit.toLayered}
\difficulty{1}
For every Boolean circuit $c$ on $n$ variables, every layer of its associated layered
feedforward circuit has a finite set of nodes.
\end{theorem}

\begin{theorem}[Gate operations in the layered circuit]
\lean{BoolCircuit.TreeCircuit.toLayered_onlyUsesGates}
\label{BoolCircuit.TreeCircuit.toLayered_onlyUsesGates}
\leanok
\uses{BoolCircuit.DAGCircuit.toLayered_onlyUsesGates, BoolCircuit.LayeredCircuit.onlyUsesGates, BoolCircuit.TreeCircuit.toLayered, BoolCircuit.stdGateOps}
\difficulty{1}
For every Boolean circuit $c$ on $n$ variables, every gate operation in its associated
layered feedforward circuit is either the identity operation, the unary operation $x
\mapsto 1-x_0$, or an AND operation of some finite arity, given by $x \mapsto \prod_i
x_i$ on inputs in $\mathrm{Fin}\,2$.
\end{theorem}

\begin{theorem}[Circuit size bound]
\lean{BoolCircuit.rewrite_size_le}
\label{BoolCircuit.rewrite_size_le}
\leanok
\difficulty{2}
Let $n,G,S,T$ be natural numbers with $S=n+G$ and $T\le n+G(S+2)$. Then $T\le S(S+2)$.
\end{theorem}

\begin{theorem}[Layered polynomial size implies polynomial size]
\lean{Language.InLayeredPPoly.inPPoly}
\label{Language.InLayeredPPoly.inPPoly}
\leanok
\uses{BoolCircuit.DAGCircuit, BoolCircuit.DAGCircuit.IsWellFormed, BoolCircuit.DAGCircuit.binarize, BoolCircuit.DAGCircuit.binarize_eval, BoolCircuit.DAGCircuit.binarize_isFaninTwo, BoolCircuit.DAGCircuit.binarize_size_le, BoolCircuit.DAGCircuit.size, BoolCircuit.DAGCircuitFamily.mem_language_iff, BoolCircuit.LayeredCircuit.toDAG, BoolCircuit.LayeredCircuit.toDAG_eval, BoolCircuit.LayeredCircuit.toDAG_isWellFormed, BoolCircuit.LayeredCircuit.toDAG_size, BoolCircuit.LayeredCircuitFamily.mem_language_iff, BoolCircuit.b2f_eq_one_iff, BoolCircuit.rewrite_size_le, Language.InLayeredPPoly, Language.InPPoly, Language.inLayeredPPoly_iff, Language.inPPoly_iff}
\difficulty{5}
If a language over Boolean strings is recognized by a family of finite layered Boolean
circuits using only the standard gate operations, with size bounded by a function $T$,
and there are natural numbers $a,k$ such that $T(n)\leq a(n+1)^k$ for every natural
number $n$, then the language has polynomial size in the sense that it is recognized by
a family of Boolean circuits with at most two inputs per gate and size at most
$a'(n+1)^{k'}$ for some natural numbers $a',k'$ at every input length $n$.
\end{theorem}

\begin{theorem}[Polynomial size in layered circuits]
\lean{Language.InPPoly.inLayeredPPoly}
\label{Language.InPPoly.inLayeredPPoly}
\leanok
\uses{BoolCircuit.DAGCircuit, BoolCircuit.DAGCircuit.IsWellFormed, BoolCircuit.DAGCircuit.deMorgan, BoolCircuit.DAGCircuit.deMorgan_eval, BoolCircuit.DAGCircuit.deMorgan_isWellFormed, BoolCircuit.DAGCircuit.deMorgan_kind_ne_or, BoolCircuit.DAGCircuit.deMorgan_size_le, BoolCircuit.DAGCircuit.size, BoolCircuit.DAGCircuit.toLayered, BoolCircuit.DAGCircuit.toLayered_eval₁, BoolCircuit.DAGCircuit.toLayered_finite, BoolCircuit.DAGCircuit.toLayered_onlyUsesGates, BoolCircuit.DAGCircuit.toLayered_size_le, BoolCircuit.DAGCircuitFamily.HasFaninTwo.isWellFormed, BoolCircuit.DAGCircuitFamily.IsWellFormed, BoolCircuit.DAGCircuitFamily.mem_language_iff, BoolCircuit.LayeredCircuit.size, BoolCircuit.LayeredCircuitFamily.mem_language_iff, BoolCircuit.b2f_eq_one_iff, BoolCircuit.rewrite_size_le, Language.InLayeredPPoly, Language.InPPoly, Language.inLayeredPPoly_iff, Language.inPPoly_iff}
\difficulty{5}
Let $L$ be a language of Boolean strings. If $L$ is recognized by a well-formed Boolean
circuit family with fan-in at most two and size at most $a(n+1)^k$ on inputs of length
$n$, for some natural numbers $a,k$, then $L$ is recognized by a finite-layered circuit
family using only identity, NOT, and unbounded fan-in AND gates, with size at most
$a'(n+1)^{k'}$ for some natural numbers $a',k'$.
\end{theorem}

\begin{theorem}[DAG and layered polynomial-size equivalence]
\lean{Language.inPPoly_iff_inLayeredPPoly}
\label{Language.inPPoly_iff_inLayeredPPoly}
\leanok
\uses{Language.InLayeredPPoly, Language.InLayeredPPoly.inPPoly, Language.InPPoly, Language.InPPoly.inLayeredPPoly}
\difficulty{1}
For every Boolean language $L$, recognition by polynomial-size families of directed
acyclic Boolean circuits with fan-in at most two is equivalent to recognition by
polynomial-size layered circuit families using only identity, NOT, and unbounded fan-in
AND gates.
\end{theorem}

\begin{theorem}[Equality of polynomial-size circuit classes]
\lean{BoolCircuit.PPoly_eq_LayeredPPoly}
\label{BoolCircuit.PPoly_eq_LayeredPPoly}
\leanok
\uses{BoolCircuit.LayeredPPoly, BoolCircuit.PPoly, Language.inPPoly_iff_inLayeredPPoly}
\difficulty{3}
The class of languages computable by families of fan-in-two DAG circuits of polynomial
size is equal to the class of languages computable by finite-layer circuit families
using the specified standard gate operations and having polynomial size.
\end{theorem}
